% unit: u34
% sources: sources/098-linear-transformations-and-their-matrices.txt, sources/059-lecture-30-linear-transformations-and-their-matrices.txt, sources/221-mit18-06scf11-ses3-6sum-pdf-download-file.txt, sources/060-problem-solving-linear-transformations.txt, sources/222-mit18-06scf11-ses3-6prob-pdf-download-file.txt, sources/223-mit18-06scf11-ses3-6sol-pdf-download-file.txt
\section{Linear transformations and their matrices}\label{sec:u34}

Many linear algebra courses begin where this unit begins: with linear transformations,
described without coordinates. A projection onto a line or a rotation of the plane makes sense
before anyone writes down a matrix, and some people (physicists especially) prefer to think
that way, about the whole space at once. But as soon as we want to compute, we choose a basis,
coordinates appear, and every linear transformation becomes a matrix. Everything we have
learned about matrices (nullspaces, determinants, eigenvalues) then applies. This unit
explains that passage from transformation to matrix, and why the eigenvector basis gives the
best matrix of all.

\subsection{Linearity without coordinates}

\begin{definition}[Linear transformation]\label{def:u34-linear}
Let $V$ and $W$ be vector spaces. A \emph{transformation} (or map) $T\colon V\to W$ assigns an
output $T(\vb{v})\in W$ to every input $\vb{v}\in V$. It is \emph{linear} if, for all vectors
$\vb{v},\vb{w}$ and all scalars $c$,
\[
  T(\vb{v}+\vb{w})=T(\vb{v})+T(\vb{w})\qquad\text{and}\qquad T(c\vb{v})=c\,T(\vb{v}).
\]
Equivalently, $T(c\vb{v}+d\vb{w})=c\,T(\vb{v})+d\,T(\vb{w})$: $T$ respects the two operations on vectors,
addition and scalar multiplication, and so respects every linear combination.
\end{definition}

\begin{proposition}[Zero goes to zero]\label{prop:u34-zero}
Every linear transformation satisfies $T(\vb{0})=\vb{0}$.
\end{proposition}
\begin{proof}
Take $c=0$ in $T(c\vb{v})=cT(\vb{v})$: $T(\vb{0})=T(0\vb{v})=0\,T(\vb{v})=\vb{0}$.
\end{proof}

\emph{Examples.}
\begin{itemize}
  \item \emph{Projection} $T\colon\R^2\to\R^2$ onto a fixed line through the origin. Doubling
        $\vb{v}$ doubles its projection, reversing $\vb{v}$ reverses it, and the projection of a sum is the
        sum of the projections.
  \item \emph{Rotation} of the plane by $45^\circ$ about the origin. Rotating $\vb{v}+\vb{w}$ gives the
        same result as rotating $\vb{v}$ and $\vb{w}$ and then adding: the whole parallelogram rotates.
  \item \emph{Multiplication by a matrix}, $T(\vb{v})=A\vb{v}$. This is linear because
        $A(\vb{v}+\vb{w})=A\vb{v}+A\vb{w}$ and $A(c\vb{v})=c\,A\vb{v}$. Any $2\times3$ matrix gives a linear
        $T\colon\R^3\to\R^2$.
  \item \emph{Differentiation}, $T(f)=f'$, on a space of functions: $(f+g)'=f'+g'$ and $(cf)'=cf'$.
        This is why we can differentiate any combination of $\sin x$, $\cos x$, $e^x$, \dots{} once we
        know the derivatives of a short list.
\end{itemize}

\emph{Non-examples.}
\begin{itemize}
  \item \emph{Shift the whole plane}, $T(\vb{v})=\vb{v}+\vb{v}_0$ with a fixed $\vb{v}_0\ne\vb{0}$. Then
        $T(2\vb{v})=2\vb{v}+\vb{v}_0$ but $2T(\vb{v})=2\vb{v}+2\vb{v}_0$, and $T(\vb{0})=\vb{v}_0\ne\vb{0}$.
  \item \emph{Length}, $T(\vb{v})=\|\vb{v}\|$ from $\R^3$ to $\R$. It does send $\vb{0}$ to $0$ and doubles
        when $\vb{v}$ doubles, but $T(-2\vb{v})=2\|\vb{v}\|\ne-2\|\vb{v}\|$.
\end{itemize}

The advantage of the coordinate-free view is the big picture. Instead of following one vector,
watch what $T$ does to a whole figure, such as the outline of a house (Figure~\ref{fig:u34-house}).
The matrix $\left[\begin{smallmatrix}1&0\\0&-1\end{smallmatrix}\right]$ keeps each $x$-component and
reverses each $y$-component, so it reflects the house across the $x$-axis and turns it upside
down; the rotation turns it without distortion; the shift moves it, origin and all, which is how
we see at once that it is not linear.

\begin{figure}[ht]\centering
\begin{tikzpicture}[>=Stealth,scale=0.85,
  house/.pic={
    \draw[thick,fill=blue!10] (0,0) -- (1,0) -- (1,1) -- (0.5,1.55) -- (0,1) -- cycle;
    \draw[thick,fill=orange!40] (0.38,0) -- (0.38,0.5) -- (0.62,0.5) -- (0.62,0);
    \fill[red!70!black] (0.5,1.55) circle (1.6pt);
  }]
  \begin{scope}
    \draw[->,gray] (-1.3,0) -- (1.6,0); \draw[->,gray] (0,-1.7) -- (0,1.9);
    \pic at (0,0) {house};
    \node[font=\small] at (0.15,-2.15) {input};
  \end{scope}
  \begin{scope}[xshift=3.9cm]
    \draw[->,gray] (-1.3,0) -- (1.6,0); \draw[->,gray] (0,-1.7) -- (0,1.9);
    \pic[rotate=45] at (0,0) {house};
    \node[font=\small] at (0.15,-2.15) {rotation by $45^\circ$};
  \end{scope}
  \begin{scope}[xshift=7.8cm]
    \draw[->,gray] (-1.3,0) -- (1.6,0); \draw[->,gray] (0,-1.7) -- (0,1.9);
    \pic[yscale=-1] at (0,0) {house};
    \node[font=\small] at (0.15,-2.15) {$\left[\begin{smallmatrix}1&0\\0&-1\end{smallmatrix}\right]$: reflection};
  \end{scope}
  \begin{scope}[xshift=11.7cm]
    \draw[->,gray] (-1.3,0) -- (1.6,0); \draw[->,gray] (0,-1.7) -- (0,1.9);
    \pic at (0.6,-1.2) {house};
    \draw[->,thick,red!70!black] (0,0) -- (0.6,-1.2) node[pos=0.45,left,font=\scriptsize] {$\vb{v}_0$};
    \node[font=\small] at (0.15,-2.15) {shift: not linear};
  \end{scope}
\end{tikzpicture}
\caption{A linear transformation acts on the whole plane at once; following a house shows the
effect. Rotation and the reflection matrix keep the origin fixed and carry straight lines to
straight lines and parallelograms to parallelograms. The shift $T(\vb{v})=\vb{v}+\vb{v}_0$ moves the corner at
the origin to $\vb{v}_0$, violating $T(\vb{0})=\vb{0}$.}
\label{fig:u34-house}
\end{figure}

\subsection{A basis determines everything}

How much must we be told about a linear $T$ to know $T(\vb{v})$ for every input? Knowing $T(\vb{v}_1)$
gives $T(c\vb{v}_1)=cT(\vb{v}_1)$ on a whole line. Knowing $T(\vb{v}_1)$ and $T(\vb{v}_2)$ for independent
$\vb{v}_1,\vb{v}_2$ gives $T(c\vb{v}_1+d\vb{v}_2)$ on a whole plane. (If $\vb{v}_2$ were $6\vb{v}_1$, then $T(\vb{v}_2)$
would carry no new information.) A basis is enough.

\begin{theorem}[Linear maps are determined by a basis]\label{thm:u34-basis}
Let $\vb{v}_1,\dots,\vb{v}_n$ be a basis of $V$. If $T$ is linear, then for every $\vb{v}=c_1\vb{v}_1+\cdots+c_n\vb{v}_n$,
\[
  T(\vb{v})=c_1T(\vb{v}_1)+c_2T(\vb{v}_2)+\cdots+c_nT(\vb{v}_n).
\]
So $T$ is completely known once $T(\vb{v}_1),\dots,T(\vb{v}_n)$ are known.
\end{theorem}
\begin{proof}
Because the $\vb{v}_i$ span $V$, every $\vb{v}$ is such a combination; linearity, applied repeatedly to
the combination, gives the formula.
\end{proof}

\begin{definition}[Coordinates]\label{def:u34-coordinates}
If $\vb{v}_1,\dots,\vb{v}_n$ is a basis of $V$ and $\vb{v}=c_1\vb{v}_1+\cdots+c_n\vb{v}_n$, the numbers
$(c_1,\dots,c_n)$ are the \emph{coordinates} of $\vb{v}$ in this basis. They are unique: if also
$\vb{v}=\sum d_i\vb{v}_i$, then $\sum(c_i-d_i)\vb{v}_i=\vb{0}$, and independence forces every $c_i=d_i$.
\end{definition}

Coordinates come from a basis. When someone writes $\vb{v}=(3,2,4)$ we silently assume the
standard basis, $\vb{v}=3\vb{e}_1+2\vb{e}_2+4\vb{e}_3$, and then the coordinates are the components. But
another basis, for example the eigenvectors of some matrix, gives the same vector different
coordinates.

\subsection{The matrix of a transformation}

To turn $T\colon V\to W$ into a matrix we choose \emph{two} bases: an input basis $\vb{v}_1,\dots,\vb{v}_n$
of $V$ and an output basis $\vb{w}_1,\dots,\vb{w}_m$ of $W$. The matrix will take the input coordinates
to the output coordinates.

\begin{definition}[Matrix of \texorpdfstring{$T$}{T} in chosen bases]\label{def:u34-matrix}
Expand the output of each input basis vector in the output basis:
\[
  T(\vb{v}_j)=a_{1j}\vb{w}_1+a_{2j}\vb{w}_2+\cdots+a_{mj}\vb{w}_m\qquad(j=1,\dots,n).
\]
The $m\times n$ matrix $A=(a_{ij})$ is the \emph{matrix of $T$} in these bases: column $j$ holds the
output coordinates of $T(\vb{v}_j)$.
\end{definition}

\begin{theorem}[The matrix does what \texorpdfstring{$T$}{T} does]\label{thm:u34-matrix}
If $\vb{v}$ has input coordinates $\vb{c}=(c_1,\dots,c_n)$, then $T(\vb{v})$ has output coordinates $A\vb{c}$.
\end{theorem}
\begin{proof}
By Theorem~\ref{thm:u34-basis} and the definition of the columns,
\[
  T(\vb{v})=\sum_{j}c_jT(\vb{v}_j)=\sum_{j}c_j\sum_{i}a_{ij}\vb{w}_i=\sum_{i}\Bigl(\sum_{j}a_{ij}c_j\Bigr)\vb{w}_i ,
\]
and $\sum_ja_{ij}c_j$ is component $i$ of $A\vb{c}$. As a sanity check, the input coordinates
$(1,0,\dots,0)$ describe $\vb{v}_1$, and $A$ times them is column~1, the coordinates of $T(\vb{v}_1)$.
\end{proof}

\begin{example}[One projection, two matrices]\label{ex:u34-projection}
Let $T$ project $\R^2$ onto the $45^\circ$ line, using the same basis for inputs and outputs.

\emph{A good basis.} Take $\vb{v}_1=(1,1)/\sqrt2$ on the line and $\vb{v}_2=(-1,1)/\sqrt2$ perpendicular to
it. The projection leaves $\vb{v}_1$ alone and kills $\vb{v}_2$: $T(\vb{v}_1)=1\vb{v}_1+0\vb{v}_2$ and
$T(\vb{v}_2)=0\vb{v}_1+0\vb{v}_2$. So
\[
  A=\begin{bmatrix}1&0\\0&0\end{bmatrix},\qquad A\begin{bmatrix}c_1\\c_2\end{bmatrix}=\begin{bmatrix}c_1\\0\end{bmatrix}.
\]
These basis vectors are the eigenvectors of the projection (eigenvalues $1$ and $0$), and the
matrix came out as the diagonal eigenvalue matrix $\Lambda$.

\emph{The standard basis.} With $\vb{v}_1=\vb{w}_1=(1,0)$ and $\vb{v}_2=\vb{w}_2=(0,1)$: the projection of $(1,0)$
onto the line is $(\tfrac12,\tfrac12)$, and so is the projection of $(0,1)$. So
\[
  P=\begin{bmatrix}\tfrac12&\tfrac12\\\tfrac12&\tfrac12\end{bmatrix}=\frac{\vb{a}\vb{a}\T}{\vb{a}\T\vb{a}}\quad\text{with }\vb{a}=(1,1),
\]
the projection matrix of Section~\ref{sec:u17}. It is symmetric with $P^2=P$, but those facts are
obvious only in the eigenvector basis, where the matrix is diagonal. The course has mostly used
the handiest basis; the best basis is the eigenvector basis.
\end{example}

\begin{example}[Differentiation as a matrix]\label{ex:u34-derivative}
Let $T=\frac{d}{dx}$ take the quadratics $c_1+c_2x+c_3x^2$ (input basis $1,x,x^2$) to the linear
polynomials (output basis $1,x$): $T(c_1+c_2x+c_3x^2)=c_2+2c_3x$. Column by column,
$T(1)=0$, $T(x)=1$, $T(x^2)=2x$, so
\[
  A=\begin{bmatrix}0&1&0\\0&0&2\end{bmatrix},\qquad
  A\begin{bmatrix}c_1\\c_2\\c_3\end{bmatrix}=\begin{bmatrix}c_2\\2c_3\end{bmatrix}.
\]
For instance $p(x)=5+3x-4x^2$ has coordinates $(5,3,-4)$, and $A(5,3,-4)=(3,-8)$, the coordinates
of $p'(x)=3-8x$. The $2\times3$ shape reflects a three-dimensional input space and a
two-dimensional output space; the nullspace of $A$, spanned by $(1,0,0)$, is the constant
polynomials, whose derivative is zero.
\end{example}

\begin{example}[The transpose as a linear map]\label{ex:u34-transpose}
On the four-dimensional space of $2\times2$ matrices let $T(X)=X\T$. It is linear:
$(X+Y)\T=X\T+Y\T$ and $(cX)\T=cX\T$. Applying it twice gives back $X$, so $T^2=I$ and $T^{-1}=T$.
\begin{itemize}
  \item In the standard basis $\vb{v}_1=\left[\begin{smallmatrix}1&0\\0&0\end{smallmatrix}\right]$,
        $\vb{v}_2=\left[\begin{smallmatrix}0&1\\0&0\end{smallmatrix}\right]$,
        $\vb{v}_3=\left[\begin{smallmatrix}0&0\\1&0\end{smallmatrix}\right]$,
        $\vb{v}_4=\left[\begin{smallmatrix}0&0\\0&1\end{smallmatrix}\right]$ we have
        $T(\vb{v}_1)=\vb{v}_1$, $T(\vb{v}_2)=\vb{v}_3$, $T(\vb{v}_3)=\vb{v}_2$, $T(\vb{v}_4)=\vb{v}_4$, so the matrix is the
        permutation
        \[
          M_T=\begin{bmatrix}1&0&0&0\\0&0&1&0\\0&1&0&0\\0&0&0&1\end{bmatrix}.
        \]
        Check: $\left[\begin{smallmatrix}1&2\\3&4\end{smallmatrix}\right]$ has coordinates $(1,2,3,4)$ and
        $M_T(1,2,3,4)=(1,3,2,4)$, the coordinates of its transpose
        $\left[\begin{smallmatrix}1&3\\2&4\end{smallmatrix}\right]$.
  \item In the basis $\vb{w}_1=\left[\begin{smallmatrix}1&0\\0&0\end{smallmatrix}\right]$,
        $\vb{w}_2=\left[\begin{smallmatrix}0&1\\1&0\end{smallmatrix}\right]$,
        $\vb{w}_3=\left[\begin{smallmatrix}0&0\\0&1\end{smallmatrix}\right]$,
        $\vb{w}_4=\left[\begin{smallmatrix}0&1\\-1&0\end{smallmatrix}\right]$, the first three are symmetric and the
        last is antisymmetric: $T(\vb{w}_i)=\vb{w}_i$ for $i\le3$ and $T(\vb{w}_4)=-\vb{w}_4$. The matrix is
        $\operatorname{diag}(1,1,1,-1)$. These $\vb{w}$'s are eigenvectors of $T$: the symmetric matrices
        are its eigenspace for $\lambda=1$ and the antisymmetric ones for $\lambda=-1$.
\end{itemize}
\end{example}

\begin{trap}
A matrix of $T$ means nothing until both bases are named, and the same $T$ has many matrices
(compare the two projection matrices above). Also, the coordinate vector of an output must be
taken in the \emph{output} basis: for $\frac{d}{dx}$, the column for $x^2$ is $(0,2)$ because
$2x=0\cdot1+2\cdot x$, not $(0,0,2)$ or $(2,0)$.
\end{trap}

\subsection{Products of transformations are products of matrices}

If $T_1\colon U\to V$ and $T_2\colon V\to W$ are linear, so is the composition $T_2\circ T_1$
($\vb{u}\mapsto T_2(T_1(\vb{u}))$). Matrix multiplication was defined to match.

\begin{proposition}[Composition and inverses]\label{prop:u34-composition}
Fix bases of $U$, $V$, $W$. If $T_1$ has matrix $A_1$ and $T_2$ has matrix $A_2$, then
$T_2\circ T_1$ has matrix $A_2A_1$. If $T\colon V\to V$ is invertible with matrix $A$ (same basis in
and out), then $T^{-1}$ has matrix $A^{-1}$.
\end{proposition}
\begin{proof}
Let $\vb{u}$ have coordinates $\vb{c}$. By Theorem~\ref{thm:u34-matrix}, $T_1(\vb{u})$ has coordinates $A_1\vb{c}$,
and then $T_2(T_1(\vb{u}))$ has coordinates $A_2(A_1\vb{c})=(A_2A_1)\vb{c}$; a matrix that does this for every
$\vb{c}$ is the matrix of the composition (its columns come from $\vb{c}=\vb{e}_j$). For the inverse,
$T^{-1}\circ T$ is the identity, whose matrix is $I$, so the matrix $B$ of $T^{-1}$ satisfies $BA=I$.
\end{proof}

\begin{example}[The second derivative]\label{ex:u34-second}
Follow $\frac{d}{dx}$ from quadratics to linear polynomials (matrix $A_1$ above) by
$\frac{d}{dx}$ from linear polynomials (basis $1,x$) to constants (basis $1$), whose matrix is
$A_2=\begin{bmatrix}0&1\end{bmatrix}$. The product
\[
  A_2A_1=\begin{bmatrix}0&1\end{bmatrix}\begin{bmatrix}0&1&0\\0&0&2\end{bmatrix}=\begin{bmatrix}0&0&2\end{bmatrix}
\]
sends $(c_1,c_2,c_3)$ to $2c_3$, and indeed $\frac{d^2}{dx^2}(c_1+c_2x+c_3x^2)=2c_3$.
\end{example}

\subsection{Exercises}

\begin{exercise}[Doubling every distance]\label{exr:u34-polar}
In polar coordinates let $T(r,\theta)=(2r,\theta)$: $T$ doubles the distance from the origin and keeps
the direction.
\begin{enumerate}[(a)]
  \item Is $T$ linear?
  \item Describe $T$ in Cartesian coordinates and check that it doubles lengths.
  \item Find its matrix in the standard basis.
\end{enumerate}
\end{exercise}
\begin{solution}
\begin{enumerate}[(a)]
  \item Yes. As a map of vectors, $T(\vb{v})=2\vb{v}$, so $T(\vb{v}_1+\vb{v}_2)=2\vb{v}_1+2\vb{v}_2=T(\vb{v}_1)+T(\vb{v}_2)$ and
        $T(c\vb{v})=2c\vb{v}=cT(\vb{v})$.
  \item $T(x,y)=(2x,2y)$, and $\|(2x,2y)\|=\sqrt{4x^2+4y^2}=2\sqrt{x^2+y^2}=2\|(x,y)\|$.
  \item $T(1,0)=(2,0)$ and $T(0,1)=(0,2)$, so the matrix is
        $\left[\begin{smallmatrix}2&0\\0&2\end{smallmatrix}\right]=2I$.
\end{enumerate}
\end{solution}

\begin{exercise}[Fixing zero is not enough]\label{exr:u34-nonlinear}
Find a transformation of $\R^2$ with $T(\vb{0})=\vb{0}$ that is not linear.
\end{exercise}
\begin{solution}
One option acts differently on different parts of the plane: $T(x,y)=(|x|,|y|)$. Then
$T\bigl((1,1)+(1,-1)\bigr)=T(2,0)=(2,0)$, but $T(1,1)+T(1,-1)=(1,1)+(1,1)=(2,2)$. Another uses a
nonlinear formula: $T(x,y)=(x,y^2)$ gives $T(c(x,y))=(cx,c^2y^2)$, which differs from
$cT(x,y)=(cx,cy^2)$ when $c\ne0,1$ and $y\ne0$. Both send $\vb{0}$ to $\vb{0}$, so $T(\vb{0})=\vb{0}$ is a
necessary condition for linearity, not a sufficient one.
\end{solution}

\begin{exercise}[Rotations compose by adding angles]\label{exr:u34-rotation}
Find the matrix $R_\theta$ of rotation by $\theta$ in the standard basis, and use
Proposition~\ref{prop:u34-composition} to show $R_\theta R_\varphi=R_{\theta+\varphi}$. What is $R_\theta^{-1}$?
\end{exercise}
\begin{solution}
Rotating $(1,0)$ gives $(\cos\theta,\sin\theta)$ and rotating $(0,1)$ gives $(-\sin\theta,\cos\theta)$, so
\[
  R_\theta=\begin{bmatrix}\cos\theta&-\sin\theta\\\sin\theta&\cos\theta\end{bmatrix}.
\]
Rotating by $\varphi$ and
then by $\theta$ is rotation by $\theta+\varphi$, so the product of the matrices must be $R_{\theta+\varphi}$.
Multiplying out confirms it: the $(1,1)$ entry is $\cos\theta\cos\varphi-\sin\theta\sin\varphi=\cos(\theta+\varphi)$ and
the $(2,1)$ entry is $\sin\theta\cos\varphi+\cos\theta\sin\varphi=\sin(\theta+\varphi)$: the matrix product contains the
addition formulas of trigonometry. The inverse rotation is $R_{-\theta}=R_\theta\T$.
\end{solution}

\begin{exercise}[Integration undoes differentiation, on one side]\label{exr:u34-integral}
Let $S$ take linear polynomials $c_1+c_2x$ to quadratics by $S(p)=\int_0^xp(t)\,dt$, and let $D=\frac{d}{dx}$
take quadratics to linear polynomials (bases $1,x$ and $1,x,x^2$). Find the matrices of $S$ and $D$, and
compute $DS$ and $SD$. Interpret.
\end{exercise}
\begin{solution}
$S(1)=x$ and $S(x)=\tfrac12x^2$, so $S$ has matrix $\left[\begin{smallmatrix}0&0\\1&0\\0&1/2\end{smallmatrix}\right]$; $D$
has matrix $\left[\begin{smallmatrix}0&1&0\\0&0&2\end{smallmatrix}\right]$ (Example~\ref{ex:u34-derivative}).
Then $DS=\left[\begin{smallmatrix}1&0\\0&1\end{smallmatrix}\right]=I$: differentiating an integral returns the
original function. But $SD=\left[\begin{smallmatrix}0&0&0\\0&1&0\\0&0&1\end{smallmatrix}\right]\ne I$: integrating a
derivative loses the constant term $c_1$. So $S$ is a right inverse of $D$ but not a left inverse,
a theme of Section~\ref{sec:u36}.
\end{solution}
